% ===================================================================
%  LENTELĖ Nr. 5 — Namas ir jo statyba.
%  Traduction 2026 d'apres texto/io, controlee sur texto/fr, texto/ru
%  et texto/cs. Consigne : voir texto/lt/10-lentele-01.tex.
%
%  LE CHANTIER EST, POUR LE GALICIEN, L'UN DES ENDROITS OU LA
%  DISTANCE A LA LANGUE VOISINE EST MINIMALE. LE LITUANIEN LE
%  CONFIRME, ET IL DIT EN PLUS POURQUOI — parce que la ligne de
%  partage n'est pas entre l'outil et le materiau, comme on
%  s'attendrait, mais entre DEUX AGES DU METIER.
%
%  Tout ce qui touche au BOIS est lituanien de fond : « dailidė » le
%  charpentier, « stalius » le menuisier, « kirvis » la hache,
%  « pjūklas » la scie, « kaltas » le ciseau, « vinis » le clou,
%  « plaktukas » le maillet, « sija » la solive, « lentos » les
%  planches. On batissait en bois ici avant qu'on y batisse
%  autrement, et la langue n'a rien eu a emprunter.
%
%  Tout ce qui touche a la PIERRE arrive avec la chose : « mūras » le
%  mur macon, « mūrininkas » le macon qui en derive, « plyta » la
%  brique, « kalkės » la chaux, « gipsas » le platre, « čerpė » la
%  tuile, « cementas », « cinkas », « skarda » la tole, « oblius » le
%  rabot. La maconnerie est venue avec l'eglise et la ville,
%  c'est-a-dire du dehors, et son vocabulaire avec elle.
%
%  C'EST LE CRITERE DU LIEU D'APPRENTISSAGE PRIS AU CHANTIER. Le
%  tableau 2 avait oppose le corps au sport, le tableau 4 le mal
%  qu'on sent au mal qu'on diagnostique ; ici s'opposent le metier
%  qu'on avait et le metier qu'on a recu. Trois pages, une seule
%  regle, et chaque fois la coupure tombe ailleurs qu'entre les
%  domaines.
%
%  DEUX GENITIFS TENUS PAR UN AUTRE TOUR. « la tri gradi (14) dil
%  perono (9) » et « la solio (10) dil eniro-pordo (8) » : ranges a
%  la lituanienne, le perron et la porte passeraient devant leurs
%  marches et leur seuil. On monte donc « trimis laipteliais (14) į
%  prieangį (9) » et l'on franchit « slenkstį (10) prie įėjimo durų
%  (8) » — deux tours aussi lituaniens que l'autre, et qui gardent
%  l'ordre de la source. Meme parti qu'au tableau 4.
%
%  LES DEUX NOMS PROPRES SONT TRADUITS, PARCE QUE L'IDO LES TRADUIT.
%  Le fac-simile francais dit Blanc et Roux ; l'ido dit Blanko et
%  Rufo, c'est-a-dire qu'il les a rendus, non transcrits. La colonne
%  tcheque avait garde les formes francaises ; on prend ici l'autre
%  parti, celui de la source qu'on traduit : « ponas Baltas » et
%  « ponas Rudys », deux noms de famille lituaniens qui existent et
%  qui disent la meme couleur. Les deux partis se defendent ; celui-ci
%  a pour lui que texto/io est le texte de depart.
%
%  ET LE LITUANIEN ENTRE DANS LA NOTE DU FEUILLET 35, comme le
%  tcheque y etait entre : la note enumere les mots nationaux qui
%  rendent le francais « solive », et « liet. sija » y prend sa
%  place. Une note qui compare des langues doit compter celle qui la
%  lit.
% ===================================================================

%%K t05-apar-1 apar
\VUsaut{6.0mm}
\VUtitre{15.0pt}{60}{LENTELĖ Nr. 5}
\VUsaut{1.4mm}
\VUtitre{11.4pt}{0}{Namas ir jo statyba.}
\VUsaut{2.2mm}

%%K t05-tit-1 sub
\VUpk{10.2pt}{40}{Laiminga sutiktis.}
\VUsaut{3.0mm}

%%K t05-01-1 p
\textsc{Jonas}. --- Dėde, labai norėčiau sužinoti, kaip statomas \VUgras{namas}.

%%K t05-01-2 p
\textsc{Dėdė} (80). --- Tai labai lengva, mano drauge, eik su manimi, parodysiu tau
tą, kurį ponas Bernardas (79) liepia statyti ir kuriame aš gyvensiu.

%%K t05-01-3 p
\textsc{Jonas}. --- O kas gi nubraižė šio namo \VUgras{planą} \textsuperscript{(76)}?

%%K t05-01-4 p
\textsc{Dėdė}. --- \VUgras{Architektas} \textsuperscript{(75)}; jis pateikė labai
aiškų brėžinį, kuris buvo patvirtintas, ir \VUgras{rangovas}
\textsuperscript{(77)} pradėjo darbus.

%%K t05-01-5 p
\textsc{Jonas}. --- O kas gi tas rangovas?

%%K t05-01-6 p
\textsc{Dėdė}. --- Tai ponas Baltas, tavo draugo Jokūbo tėvas. Jis vadovauja
darbininkams kartu su ponu Rudžiu, architektu.

%%K t05-01-7 p
Na štai! Kaip tik laiku ateina ponas Baltas su savo \VUgras{liniuote}
\textsuperscript{(78)} ir ponas Rudys su savo planu.

%%K t05-01-8 p
Labas rytas, ponai. Kaip laikotės?

%%K t05-01-9 p
\textsc{Ponas Rudys}. --- Labai gerai, ačiū. Labas rytas, Jonai, kaip užaugo tas
vaikas!

%%K t05-01-10 p
\textsc{Dėdė}. --- Taip, iš tikrųjų, jis mažas vyras. Jis labai rimtas, reikia jam
paaiškinti viską, ką jis mato. Šiandien jis norėtų sužinoti, kaip statomas namas.

%%K t05-tit-2 sub
\VUpk{10.2pt}{40}{Darbininkai.}
\VUsaut{3.0mm}

%%K t05-01-11 p
\textsc{Ponas Baltas}. --- Na, eik su manimi, drauguži, tuoj tau tai paaiškinsiu,
nes labai mėgstu vaikus, kurie nori mokytis.

%%K t05-01-12 p
Matai tą darbininką, kuris laiko rankoje \VUgras{kastuvą} \textsuperscript{(82)}: jis
\VUgras{žemkasys} \textsuperscript{(81)}. Jis kasa \VUgras{tranšėją}
\textsuperscript{(46)} vandentiekio vamzdžiams pakloti. Jis taip pat iškasė žemę toje
vietoje, kur stovi namas, kad mūrininkas galėtų iškelti keturias \VUgras{sienas}.
Ta sienų dalis, kuri yra žemėje, vadinama \VUgras{pamatais} \textsuperscript{(2)}.
Erdvė, apimta pamatų, sudaro \VUgras{rūsį} \textsuperscript{(3)}. Tas rūsys turi mažą
langelį, vadinamą \VUgras{rūsio langu} \textsuperscript{(5)}, \VUgras{duris}
\textsuperscript{(4)} ir laiptus.

%%K t05-01-13 p
\VUgras{Mūrininkas} \textsuperscript{(108)} toliau statė sienas, palikdamas angas
\VUgras{durims} \textsuperscript{(8)} ir \VUgras{langams} \textsuperscript{(17)}.

%%K t05-01-14 p
\textsc{Jonas}. --- Bet to mūrininko aš nematau!

%%K t05-01-15 p
P. \textsc{Baltas}. --- Dabar jis baigė dirbti name. Jis prie \VUgras{arklidės}
\textsuperscript{(54)}, ant savo \VUgras{pastolių} \textsuperscript{(110)}, stato
\VUgras{skalbyklos} \textsuperscript{(56)} sieną.

%%K t05-01-16 p
Jis turi mentelę, kalkinę ir \VUgras{svambalą} \textsuperscript{(109)}. Bet tu tų
vardų neatsiminsi.

%%K t05-01-17 p
\textsc{Jonas}. --- Tikrai atsiminsiu, nes tuoj paprašysiu dėdės mažos mentelės ir
kalkinės, o svambalą pats pasidarysiu iš virvutės.

%%K t05-tit-3 sub
\VUsaut{3.0mm}
\VUpk{10.2pt}{40}{Medžiagos.}
\VUsaut{3.0mm}

%%K t05-01-18 p
\textsc{Dėdė}. --- A! labai gerai. Bet pasakyk man, Jonai, ko reikia namui statyti?

%%K t05-01-19 p
\textsc{Jonas}. --- Reikia \VUgras{tašytų akmenų} \textsuperscript{(59)} ir
\VUgras{skiedinio} \textsuperscript{(65)}.

%%K t05-01-20 p
\textsc{Dėdė}. --- Teisingai, žiūrėk į tą vyrą su didele skrybėle ant savo
\VUgras{vežimo} \textsuperscript{(136)}. Jis \VUgras{vežikas} \textsuperscript{(135)}.
Kasdien jis atveža čia \VUgras{akmens luitų} \textsuperscript{(60)}. Jis iškrauna
savo vežimą kieme, o \VUgras{akmenskaldžiai} \textsuperscript{(83)} tašo luitus ir
pjauna juos savo \VUgras{akmenų pjūklu} \textsuperscript{(85)}. \VUgras{Juodadarbis}
\textsuperscript{(146)} neša juos mūrininkui, kuris juos guldo.

%%K t05-01-21 p
Tuo tarpu \VUgras{maišytojas} \textsuperscript{(87)} ruošia skiedinį. Jam reikia
\VUgras{smėlio} \textsuperscript{(62)}, \VUgras{kalkių} \textsuperscript{(63)} ir
\VUgras{vandens} \textsuperscript{(64)}. Kalkės netoli jo \VUgras{maiše}
\textsuperscript{(71)}, \VUgras{mūrininko padėjėjas} \textsuperscript{(111)} atveža jam
smėlio \VUgras{karučiais} \textsuperscript{(112)}, o \VUgras{pameistrys}
\textsuperscript{(113)} yra prie siurblio (58). Jis pripildo \VUgras{kibirą}
\textsuperscript{(114)}. Skiedinys netrukus paruoštas. \VUgras{Skiedinio nešėjas}
\textsuperscript{(107)} paima jį ir kopėčiomis neša aukštyn ant pastolių, kur dirba
mūrininkas, baigiantis tą namo pusę.

%%K t05-01-22 p
O dabar, kaip vadinamas darbininkas, kuris dirba \VUgras{stogą}
\textsuperscript{(31)}?

%%K t05-01-23 p
\textsc{Jonas}. --- Jis \VUgras{stogdengys} \textsuperscript{(118)}.

%%K t05-tit-4 sub
\VUsaut{3.0mm}
\VUpk{10.2pt}{40}{Namo stogas.}
\VUsaut{3.0mm}

%%K t05-01-24 p
P. \textsc{Baltas}. --- Taip, bet pirmiausia ateina \VUgras{dailidė}
\textsuperscript{(115)}.

%%K t05-01-25 p
Dailidė dirba \VUgras{stogo konstrukciją} \textsuperscript{(35)}. Jis surenka
\VUgras{sijas} \textsuperscript{(37)} \VUgras{kaiščiais} \textsuperscript{(117)}. Tie
darbininkai ten viršuje, su savo plačiomis aksominėmis rankovėmis, yra dailidės.
Vienas laiko rąstelį, o kitas pjauna kaištį savo \VUgras{kirviu}
\textsuperscript{(116)}. Tas, kuris yra prie \VUgras{kamino} \textsuperscript{(41)},
yra stogdengys. Jis kala grebėstus ant (*) sijų, o \VUgras{skalūnus}
\textsuperscript{(66)} ant grebėstų. Kartais jis deda čerpes.

%%K t05-noto-1 noto
\VUnotes{143.2mm}{%
(*) \VUgras{Balk-o} = vok. \textit{Balken}, angl. \textit{joist}, pranc.
\textit{solive}, it. \textit{travicello}, rus. \textit{balka}, liet.
\textit{sija}, isp. ir port. \textit{viga}. Techninė šaknis iki šiol nepriimta, bet
siūlytina prancūzų \textit{solive} prasmei perteikti; tai nėra nei \textit{trabo}
(pranc. \textit{poutre}), nei rąstelis. Tai aptašyta medinė kartis, kuri laiko grindis
ir lubas.%
}

%%K t05-01-26 p
Arklidės stogas \VUgras{dengtas čerpėmis} \textsuperscript{(55)}, namo ---
\VUgras{dengtas skalūnu} \textsuperscript{(32)}, o verandos --- \VUgras{cinkinis}
\textsuperscript{(24)}.

%%K t05-01-27 p
\textsc{Jonas}. --- Ar stogdengys dirba ir cinkinius stogus?

%%K t05-01-28 p
P. \textsc{Baltas}. --- Ne, bet \VUgras{cinkuotojas} \textsuperscript{(121)} arba
\VUgras{skardininkas,} tas, kuris deda \VUgras{stoglangį} \textsuperscript{(38)} ant
namo stogo. Jis deda ir \VUgras{lietvamzdžius} \textsuperscript{(43)} bei
\VUgras{latakus} \textsuperscript{(42)}. Jis lituoja cinką ir skardą. Jis pritvirtino
\VUgras{žaibolaidį} \textsuperscript{(40)} ir \VUgras{vėtrungę} \textsuperscript{(39)}
ant \VUgras{skydo} \textsuperscript{(36)}.

%%K t05-01-29 p
\textsc{Jonas}. --- Vadinasi, namas baigtas?

%%K t05-tit-5 sub
\VUsaut{3.0mm}
\VUpk{10.2pt}{40}{Įrankiai.}
\VUsaut{3.0mm}

%%K t05-01-30 p
P. \textsc{Baltas}. --- Ne visiškai. \VUgras{Tinkuotojas} \textsuperscript{(99)}
stato iš \VUgras{plytų} \textsuperscript{(61)} pertvaras, kurios dalija jį į
įvairius kambarius. Jis tinkuoja \VUgras{gipsu} \textsuperscript{(70)}
\VUgras{lubas} \textsuperscript{(26)} ir sienas. Dabar jis dirba \VUgras{verandos}
\textsuperscript{(33)} lubas. Jis turi, kaip ir mūrininkas, \VUgras{mentelę}
\textsuperscript{(100)} ir \VUgras{kalkinę} \textsuperscript{(101)}.

%%K t05-01-31 p
Paskui stalius dirba duris, langus, \VUgras{parketą} \textsuperscript{(16)} ir
\VUgras{langines} \textsuperscript{(19)}.

%%K t05-01-32 p
Dabar jis dirba kieme prie savo \VUgras{varstoto} \textsuperscript{(90)}. Jis turi
\VUgras{pjūklą} \textsuperscript{(94)} medienai (69) pjauti, \VUgras{oblių}
\textsuperscript{(91)}, \VUgras{spaustuvą} \textsuperscript{(92)}, \VUgras{kaltą}
\textsuperscript{(94 bis)}, \VUgras{skriestuvą} \textsuperscript{(95)} ir medinį
\VUgras{plaktuką} \textsuperscript{(95 bis)}.

%%K t05-01-33 p
\textsc{Jonas}. --- A! bet linksmiau dailidžiauti negu mūryti. Dėde, nupirksi man
varstotą, pjūklą ir oblių, nes netrukus dirbsiu būdą Medorui.

%%K t05-01-34 p
\textsc{Dėdė}. --- Taip, vaikeli.

%%K t05-01-35 p
\textsc{Jonas}. --- Koks aš patenkintas! O kas tas, kuris stovi prie įėjimo durų?

%%K t05-tit-6 sub
\VUsaut{3.0mm}
\VUpk{10.2pt}{40}{Dažymas.}
\VUsaut{3.0mm}

%%K t05-01-36 p
P. \textsc{Baltas}. --- Jis \VUgras{dažytojas} \textsuperscript{(102)}. Jis stovi ant
savo kopėčių su \VUgras{dažų} \textsuperscript{(103)} puodu ir savo \VUgras{teptuku}
\textsuperscript{(104)}. Jis dažo įėjimo durų medį, kad ilgiau išsilaikytų.

%%K t05-01-37 p
Jis taip pat klijuoja tapetus butuose.

%%K t05-01-38 p
\VUgras{Apmušėjas} \textsuperscript{(107)}, kuris yra ant \VUgras{kopėčių}
\textsuperscript{(106)}, ant \VUgras{balkono} \textsuperscript{(28)}, deda
\VUgras{užuolaidą} \textsuperscript{(29)}.

%%K t05-01-39 p
Darbininkas, stovintis prie didelio lango, yra \VUgras{stiklius}
\textsuperscript{(96)}. Jis tvirtina \VUgras{stiklus} \textsuperscript{(18)}
\VUgras{glaistu} \textsuperscript{(73)}. Tas kitas darbininkas, atsiklaupęs prie rūsio
durų, yra \VUgras{šaltkalvis} \textsuperscript{(97)}. Jis deda \VUgras{spyną}
\textsuperscript{(6)}. Jo \VUgras{dildė} \textsuperscript{(98)} netoli jo.

%%K t05-01-40 p
\textsc{Jonas}. --- Ar irgi šaltkalvis tas darbininkas, kuris savo \VUgras{plaktuku}
\textsuperscript{(84)} kala geležį?

%%K t05-01-41 p
P. \textsc{Baltas}. --- Tiksliau sakant, jis kalvis (125). Jis kala \VUgras{geležis}
\textsuperscript{(67)} \VUgras{elektrikui} \textsuperscript{(124)}, kuris yra ant
\VUgras{vežiminės} \textsuperscript{(51)} stogo. Jis turi \VUgras{žaizdrą}
\textsuperscript{(126)}, \VUgras{priekalą} \textsuperscript{(127)} ir
\VUgras{žnyples} \textsuperscript{(128)}.

%%K t05-01-42 p
Elektrikas tvirtina telefono \VUgras{varinę} \textsuperscript{(44)} \VUgras{vielą}.

%%K t05-tit-7 sub
\VUpk{10.2pt}{40}{Sodininkystė.}
\VUsaut{3.0mm}

%%K t05-01-43 p
\textsc{Dėdė}. --- \VUgras{Kiemo} \textsuperscript{(45)} gale, prie \VUgras{grotų}
\textsuperscript{(47)} ir \VUgras{vartų} \textsuperscript{(48)}, matai
\VUgras{sodininką} \textsuperscript{(129)}. Jis ruošia \VUgras{daržiuką}
\textsuperscript{(49)}. Jis perkasė žemę savo \VUgras{kastuvu}
\textsuperscript{(131)}, sėja sėklas ir lygina \VUgras{grėbliu}
\textsuperscript{(130)}. Paskui savo \VUgras{laistytuvu} \textsuperscript{(132)} jis
palies \VUgras{lysves} \textsuperscript{(150)}, ir augalai augs.

%%K t05-01-44 p
\VUgras{Arklininkas} \textsuperscript{(133)}, kuris ką tik nužiūrėjo arklį ir davė
jam pašaro, valo \VUgras{ekipažą} \textsuperscript{(52)} kempine, \VUgras{rankiniu
siurbliu} \textsuperscript{(134)} ir vandens kibirėliu. Paskui jis vėl įves jį į
vežiminę ir uždarys duris storu \VUgras{skląsčiu} \textsuperscript{(53)}.

%%K t05-01-45 p
Štai tu patenkintas, manau, gavai pakankamai paaiškinimų.

%%K t05-01-46 p
\textsc{Jonas}. --- Taip, dėde, bet noriai pamatyčiau namo vidų.

%%K t05-01-47 p
\textsc{Dėdė}. --- Tai bus rytoj. Ponas Rudys paaiškins tau planą ir nurodys tau
kiekvieno kambario vardą.

%%K t05-tit-8 sub
\VUsaut{3.0mm}
\VUpk{10.2pt}{40}{Vidinis namo išplanavimas.}
\VUsaut{2.4mm}

%%K t05-apar-2 apar
{\centering\textit{(Žiūrėk planą.)}\par}
\VUsaut{2.4mm}

%%K t05-01-48 p
\textsc{Architektas}. --- Eime, Jonai, tuoj leisiu tau apžiūrėti įvairias namo
dalis.

%%K t05-01-49 p
Įeikime į \VUgras{pirmąjį aukštą} \textsuperscript{(7)}. Štai \VUgras{skambučio}
\textsuperscript{(11)} mygtukas. Užlipame trimis \VUgras{laipteliais}
\textsuperscript{(14)} į \VUgras{prieangį} \textsuperscript{(9)}, peržengiame
\VUgras{slenkstį} \textsuperscript{(10)} prie įėjimo \VUgras{durų}
\textsuperscript{(8)}, ir štai esame prie \VUgras{laiptų} \textsuperscript{(13)}
apačios.

%%K t05-01-50 p
\textsc{Jonas}. --- Tai koridorius.

%%K t05-01-51 p
\textsc{Architektas}. --- Ne, tai \VUgras{prieškambaris} \textsuperscript{(12)}.
\VUgras{Koridorius} \textsuperscript{(a)} yra gale. Dešinėje štai \VUgras{svetainė}
\textsuperscript{(b)} ir \VUgras{valgomasis} \textsuperscript{(c)}, priešais
valgomąjį yra \VUgras{virtuvė} \textsuperscript{(d)} ir \VUgras{tarnų kambarys}
\textsuperscript{(e)}, o toliau priekyje \VUgras{darbo kambarys}
\textsuperscript{(f)}. \VUgras{Veranda} \textsuperscript{(23)} su savo
\VUgras{stiklinėmis durimis} \textsuperscript{(25)} yra prie svetainės.

%%K t05-01-52 p
Tuoj užlipsime laiptais. Laikykis už \VUgras{turėklo} \textsuperscript{(15)} ir
skaičiuok laiptelius. Jų keturiolika. Štai \VUgras{aikštelė}
\textsuperscript{(m)}, esame antrajame \VUgras{aukšte} \textsuperscript{(27)}.

%%K t05-tit-9 sub
\VUsaut{3.0mm}
\VUpk{10.2pt}{40}{Antrasis aukštas.}
\VUsaut{3.0mm}

%%K t05-01-53 p
Priešais laiptus yra \VUgras{išvietė} \textsuperscript{(20)} ir \VUgras{prausykla}
\textsuperscript{(j)}. Dešinėje graži \VUgras{miegamoji} \textsuperscript{(g)} su
dviem langais į namo \VUgras{fasadą} \textsuperscript{(1)}. Kairėje yra
\VUgras{vonia} \textsuperscript{(1)} ir \VUgras{svečių kambarys}
\textsuperscript{(i)} su didele \VUgras{alkova} \textsuperscript{(h)}: prie
miegamosios yra \VUgras{vaikų kambarys} \textsuperscript{(k)}. Jis prieina prie
plačios \VUgras{terasos} \textsuperscript{(21)}. Po ta terasa yra kita išvietė (20),
o prie sienos štai \VUgras{varpas} \textsuperscript{(22)}, kuris skamba vakarienei.

%%K t05-01-54 p
\textsc{Jonas}. --- Eikime į \VUgras{trečiąjį aukštą}.

%%K t05-01-55 p
\textsc{Architektas}. --- Trečiojo aukšto nėra. Po stogu yra tik
\VUgras{mansardos} \textsuperscript{(33)} ir palėpė (34). Baigėme apžiūrą, galime
leistis atgal.

%%K t05-01-56 p
\textsc{Jonas}. --- Ačiū, pone, tai labai įdomu. Tuoj papasakosiu tėvui viską, ką
mačiau.

\VUsaut{4.0mm}
\VUfilet{22mm}
